%% sections/zhaoyan.tex: Zhaoyan's slides (post-training Qwen3-8B for catalogue decisions). OWNER: Zhaoyan.
%% Only Zhaoyan edits this file. Drafted by Renzhi on 2026-11-27 from Zhaoyan's report section (final report,
%% "Post-training Qwen3-8B for catalogue decisions", branch doc/qwen-posttrain-results); every number is from there.
%% - Figures go in figures/zhaoyan/ and are included as \includegraphics[width=...]{zhaoyan/<file>}.
%% - Need a macro of your own? Define it here with a zhaoyan prefix (\newcommand{\zhaoyanFoo}{...}), never in preamble.tex.
%% - Build just your part: latexmk -pdf solo_zhaoyan.tex   (or compile solo_zhaoyan.tex in Overleaf).
%% - Backup slides: put them at the end of this file as \begin{frame}[noframenumbering]{Backup: ...}.
%% - The section's title and cover slide are set in main.tex; ask Renzhi to rename it.

\tikzset{
  zhaoyanBox/.style={draw=usccardinal, fill=usccardinal!6, rounded corners=2pt, thick, align=center,
                     font=\footnotesize, text width=24mm, minimum height=11mm, inner sep=2pt},
  zhaoyanSim/.style={zhaoyanBox, draw=uscgray, fill=uscgray!8},
  zhaoyanArr/.style={-{Stealth[length=2mm]}, thick, uscgray},
}

\begin{frame}{The task: a small model that picks circuits}
  \centering
  \begin{tikzpicture}
    \node[zhaoyanSim] (in)  at (0,0)      {\textbf{Spec}\\ catalogue, history};
    \node[zhaoyanBox] (llm) at (34mm,0)   {\textbf{Qwen3-8B}\\ proposes 3 circuits};
    \node[zhaoyanSim] (srch) at (68mm,0)  {\textbf{Search} each:\\ 128, then up to 200 more};
    \node[zhaoyanBox] (dec) at (102mm,0)  {\textbf{Decide}\\ accept, retry, reject};
    \draw[zhaoyanArr] (in) -- (llm); \draw[zhaoyanArr] (llm) -- (srch); \draw[zhaoyanArr] (srch) -- (dec);
    \draw[zhaoyanArr, usccardinal] (dec.south) -- ++(0,-6mm) -| node[below, pos=0.25, font=\scriptsize]
      {another round (at most 3), with the measured results} (llm.south);
  \end{tikzpicture}

  \vspace{6mm}
  \begin{itemize}\small
    \item The outer loop of the designer: given a spec, choose which \alert{existing catalogue circuits} to try, read
      the search results, and decide when to stop.
    \item The search replays stored SKY130 designs: no new simulation.
    \item An accept counts only with a verified design that meets the spec.
  \end{itemize}
\end{frame}

\begin{frame}{Data and training}
  \begin{columns}[T]
    \begin{column}{0.48\textwidth}
      \textbf{Data, from stored searches}
      \small
      \begin{itemize}
        \item 6,000 successful sessions replayed from 478 training tasks; no teacher model.
        \item Successful sessions after 0, 1 or 2 failed rounds: 30\%, 50\%, 20\%.
        \item Failed rounds and their results are context; the model learns the successful proposal and the accept
          (11,977 examples).
      \end{itemize}
    \end{column}
    \begin{column}{0.48\textwidth}
      \textbf{Two stages}
      \small
      \begin{enumerate}
        \item \textbf{SFT:} LoRA (rank 32) on Qwen3-8B, 2 epochs.
        \item \textbf{GRPO:} 578 tasks $\times$ 16 complete sessions each, compared within a task.\\[2pt]
          Reward: $+1$ for a verified accept, $0$ for a reject, $-1$ for an unsupported accept or an exhausted
          session, $-0.1$ per invalid reply (at most 2).
      \end{enumerate}
      \vspace{1mm}
      \footnotesize Both checkpoints fixed before the test.
    \end{column}
  \end{columns}
\end{frame}

\begin{frame}{Results on the test split}
  \begin{columns}[c]
    \begin{column}{0.58\textwidth}
      \centering
      \includegraphics[width=\linewidth]{zhaoyan/qwen_test276_success.png}
    \end{column}
    \begin{column}{0.40\textwidth}
      \footnotesize
      276 test sessions (138 tasks $\times$ 2 seeds). Correct: a verified accept when the catalogue holds a solution
      (240), a reject when it does not (36).
      \small
      \begin{itemize}
        \item Untrained, format enforced: 74.6\%.
        \item SFT + GRPO: \alert{83.3\%}, +8.7 points (95\% interval 2.2 to 15.2).
        \item Rule designer: \alert{87.7\%}; GRPO is 4.3 points below ($-8.7$ to $-0.4$).
      \end{itemize}
    \end{column}
  \end{columns}
\end{frame}

\begin{frame}{What each stage teaches}
  \begin{columns}[c]
    \begin{column}{0.53\textwidth}
      \centering\footnotesize
      \begin{tabular}{@{}lccc@{}}
        \toprule
        276 sessions & Base & SFT & +\,GRPO \\
        \midrule
        Verified accepts (of 240) & 171 & 196 & 194 \\
        Correct rejects (of 36) & 35 & \alert{1} & 36 \\
        Invalid replies & 27 & \alert{150} & 0 \\
        \bottomrule
      \end{tabular}
    \end{column}
    \begin{column}{0.45\textwidth}
      \small
      \begin{itemize}
        \item \textbf{SFT finds more solutions} (171 $\to$ 196) but stops wrongly: after the last round it keeps
          proposing, so it rejects only 1 of the 36 sessions with no solution. Its examples all end in an accept.
        \item \textbf{GRPO fixes the stopping:} all 36 rejects right, no invalid replies, and 194 of SFT's 196
          solutions kept.
        \item Over SFT, GRPO gains by stopping correctly, not by finding more.
      \end{itemize}
    \end{column}
  \end{columns}
\end{frame}

\begin{frame}{Takeaways}
  \begin{itemize}
    \item SFT teaches the model to find solutions; GRPO teaches it when to stop. Only the two together beat the
      untrained model: +8.7 points.
    \item It still does not beat the fixed rule designer (83.3\% against 87.7\%): as for the RL methods, a simple
      rule is the bar.
    \item Compensated (two-stage) circuits: SFT and GRPO accept none, so their stability is not tested.
    \item \textbf{Next:} training data that ends at the round limit and with justified rejects; rewards for verified
      progress and checked stability; then connecting the model to the compiler and simulator and testing on new
      circuit families.
  \end{itemize}
  \vspace{2mm}
  {\footnotesize\color{uscgray} Adapter: \texttt{JzyRay/Qwen3-8B-EE599-SFT2-GRPO578} (Hugging Face); code in
  \texttt{llm/experiments/qwen\_posttrain/}.}
\end{frame}
